\documentclass[12pt,a4paper]{article}
%
% AUTHOR: Prof. RUI SANTOS CRUZ
%
\usepackage{helvet} 
\renewcommand{\familydefault}{\sfdefault}
\usepackage{a4wide}
\usepackage{ucs}
\usepackage[utf8x]{inputenc}
%%%%%%%%%%%%%%%%%%%%%%%%%%%%%%%%%%%%%%%%%%%%%%%%%%%%%%%%%%%%%%%%%%%%%%%%%%%%%%%%%%%
% SELECT ONE OF THE FOLLOWING PACKAGES FOR THE LANGUAGE 
\usepackage[english]{babel}
% \usepackage[portuges]{babel}

%%%%%%%%%%%%%%%%%%%%%%%%%%%%%%%%%%%%%%%%%%%%%%%%%%%%%%%%%%%%%%%%%%%%%%%%%%%%%%%%%%%
\usepackage{graphicx}
\usepackage{hyperref}
\usepackage{cite}
\usepackage[absolute]{textpos}
\usepackage{tabularx} 
\usepackage{tabulary}                 
\usepackage{fancyhdr}
\usepackage[table]{xcolor}
\pagestyle{fancy}
\headsep=50pt
\setlength{\headheight}{50pt}
\usepackage{listings}
\lstset{escapeinside={<@}{@>}}
\usepackage{minted}
\definecolor{LightGray}{rgb}{0.95, 0.95, 0.95}
\definecolor{darkblue}{rgb}{0.0,0.0,0.6}
\definecolor{editorOcher}{rgb}{1, 0.5, 0}

% Clever Referencing of document parts
\usepackage{cleveref}

\lstdefinestyle{commandline} {%
language={[WinXP]command.com},
breaklines=true,
%aboveskip=\baselineskip,
belowskip=\baselineskip,
showstringspaces=false,
backgroundcolor=\color{black},
basicstyle=\small\color{white}\ttfamily
showstringspaces=false,
keywordstyle=\color{cyan}\bfseries,
stringstyle=\color{gray}\ttfamily,
commentstyle=\color{green}\itshape,
moredelim=[s][\color{yellow}\bfseries]{C:}{\>}
}

\lstdefinestyle{Bash} {%
language=bash,
breaklines=true,
belowskip=\baselineskip,
backgroundcolor=\color{LightGray},
showstringspaces=false,
keywordstyle=\color{black}\bfseries,
basicstyle=\small\color{black}\ttfamily,
stringstyle=\color{editorOcher}\ttfamily,
commentstyle=\color{brown}\itshape,
otherkeywords={xcode-select, mkdir,rm},
moredelim=[s][\color{darkblue}]{~$},
literate={~} {$\sim$}{1}
}

\lstdefinestyle{Rubytext} {%
language=Ruby,
breaklines=true,
belowskip=\baselineskip,
basicstyle=\small\ttfamily\color{black},
backgroundcolor=\color{LightGray},
showstringspaces=false,
commentstyle = \ttfamily\color{red},
keywordstyle=\ttfamily\color{blue},
stringstyle=\color{orange}
}


\lstdefinestyle{py} {%
	language=python,
	literate=%
	*{0}{{{\color{red}0}}}1
	{1}{{{\color{red}1}}}1
	{2}{{{\color{red}2}}}1
	{3}{{{\color{red}3}}}1
	{4}{{{\color{red}4}}}1
	{5}{{{\color{red}5}}}1
	{6}{{{\color{red}6}}}1
	{7}{{{\color{red}7}}}1
	{8}{{{\color{red}8}}}1
	{9}{{{\color{red}9}}}1,
	basicstyle=\small\ttfamily,
	numbers=left,
	% numberstyle=\tiny,
	% stepnumber=2,
	numbersep=5pt,
	tabsize=4,
	extendedchars=true,
	breaklines=true,
	keywordstyle=\color{blue}\bfseries,
	frame=b,
	commentstyle=\color{brown}\itshape,
	stringstyle=\color{editorOcher}\ttfamily,
	showspaces=false,
	showtabs=false,
	xleftmargin=17pt,
	framexleftmargin=17pt,
	framexrightmargin=5pt,
	framexbottommargin=4pt,
	backgroundcolor=\color{LightGray},
	showstringspaces=false,
}

%%%%%%%%%%%%%%%%%%%%%%%%%%%%%%%%%%%%%%%%%%%%%%%%%%%%%%%%%%%%%%%%%%%%%%%%%%%%%%%%%%%
% PLEASE FILL THE ADEQUATE DATA IN THE TABLE REPLACING
% THE VALUES EXEMPLIFIED
\lhead{}
\renewcommand{\arraystretch}{1.1}
\fancyhead[C]{\begin{tabularx}{1.0\textwidth}{|l|X|l|l|}
\hline 
% In the following line change Course Name: PPIII, PPB
\textbf{AGISIT 19/20} & \textbf{ASSIGNMENT} & \textbf{Number:}  &  1 \\ 
\hline
% In the following line insert your Name and IST ID
\multicolumn{2}{|l|}{Assignment type (Preparation, Reading, Lab)} & \textbf{Issue Date:}  &  \today \\ 
\hline
% In the following line insert the Activity CODE and Title (abridged)
%\textbf{WP n.} (99) & (Subject) & \textbf{Group:} & (99) \\
\multicolumn{2}{|l|}{Title of Assignment} & \textbf{Due Date:} & \today \\ 
\hline
\multicolumn{2}{|l|}{Author:  } & \textbf{Revision:} & 0.0 \\ 
\hline
\end{tabularx}}
\rhead{}

%%%%%%%%%%%%%%%%%%%%%%%%%%%%%%%%%%%%%%%%%%%%%%%%%%%%%%%%%%%%%%%%%%%%%%%%%%%%%%%%%%%
% DO NOT CHANGE THIS BLOCK
\begin{document}
\textblockorigin{-18pt}{-2pt}
\begin{textblock*}{10cm}(2cm,1cm)
\includegraphics[width=6cm]{IST_C_RGB_POS-01-CROP.jpg}
\end{textblock*}
%%%%%%%%%%%%%%%%%%%%%%%%%%%%%%%%%%%%%%%%%%%%%%%%%%%%%%%%%%%%%%%%%%%%%%%%%%%%%%%%%%%
%%%%%%%%%%%%%%%%%%%%%%%%%%%%%%%%%%%%%%%%%%%%%%%%%%%%%%%%%%%%%%%%%%%%%%%%%%%%%%%%%%%
\section{Introduction}
The objective of this Lab assignment is to study and to gain a better understanding of something.




\begin{figure}[h]
\centering
\includegraphics[width=0.7\textwidth]{Ansible_env_2.pdf}
\caption{The Webfront Environment}
\label{fig:ansible-webfront}
\end{figure}

%%%%%%%%%%%%%%%%%%%%%%%%%%%%%%%%%%%%%%%%%%%%%%%%%%%%%%%%%%%%%%%%%%%%%%%%%%%%%%%%%%%
% YOUR TEXT STARTS HERE
%%%%%%%%%%%%%%%%%%%%%%%%%%%%%%%%%%%%%%%%%%%%%%%%%%%%%%%%%%%%%%%%%%%%%%%%%%%%%%%%%%%   
\section{How to do it?}
One-line code formatting works with minted. For instance, a simple html sample like this:
\mint{html}|<h2>Something <b>here</b></h2>| 
can be properly 
formatted.

Bla bla.\\

\begin{minted}[frame=lines,framesep=2mm,baselinestretch=1.2,bgcolor=LightGray]{bash}

#!/bin/sh
# removes backup files and TeX files

if [ $# -gt 0 ] && [ $1 = "-b" ]
then
    shift
    backups=false
else
    backups=true
fi

...
\end{minted}

\begin{minted}[frame=lines,framesep=2mm,baselinestretch=1.2,bgcolor=LightGray]{bash}

$ cd ~/containers/vagrant
$ mkdir jenkins-docker && cd jenkins-docker
$ vagrant init ubuntu/trusty64
$ vagrant up
\end{minted}
%%%%%%%%%%%%%%%%%%%%%%%%%%%%%%%%%%%%%%%%%%%%%%%%%%%%%%%%%%%%%%%%%%%%%%%%%%%%%%%%%%%
% YOUR TEXT ENDS HERE
%%%%%%%%%%%%%%%%%%%%%%%%%%%%%%%%%%%%%%%%%%%%%%%%%%%%%%%%%%%%%%%%%%%%%%%%%%%%%%%%%%% 
\end{document}                             % The required last line